\section{ROUND-2 OBJECTIVE}

I pursue the ``strict advancement without overclaiming'' route. The global statement is still open unconditionally, while the currently available literature gives a conditional disproof under a qualitative prime tuples conjecture. The most promising unconditional Round-2 target is therefore to prove the conjectured inequality for a substantially larger infinite family than in Round-1.

The new theorem proved in this round is the following:

\medskip

\noindent\textbf{Theorem.}
\emph{Let \(n=r^{a}s\), where \(a\geq 1\) and \(r\neq s\) are primes. Let}
\[
1=d_1<d_2<\cdots<d_t=n
\]
\emph{be the positive divisors of \(n\) in increasing order. Then}
\[
\sum_{1\leq i<j\leq t}\frac{1}{d_j-d_i}
\leq 9\left(1+\sum_{1\leq i<t}\frac{1}{d_{i+1}-d_i}\right).
\]

This strictly strengthens the Round-1 progress, which only handled prime powers.

\section{Round-1 FOUNDATION USED}

I use the following Round-1 material as fixed input.

\begin{enumerate}
\item The formal notation
\[
S_{\mathrm{all}}(n):=\sum_{1\le i<j\le t}\frac{1}{d_j-d_i},
\qquad
S_{\mathrm{gap}}(n):=\sum_{1\le i<t}\frac{1}{d_{i+1}-d_i}.
\]
\item The Round-1 conclusion that the unconditional global problem remains open, while Tao's 2025 note gives only a conditional disproof under a qualitative Hardy--Littlewood prime tuples conjecture.
\item The Round-1 prime-power result, used here only as a baseline comparison: prime powers satisfy the conjectured inequality with an absolute constant.
\end{enumerate}

No Round-1 proof is re-used verbatim below; the new argument is independent.

\section{NEW INSIGHT / TOOL (ROUND-2)}

The new tool is a structural representation of the divisor set of \(n=r^a s\).

Write
\[
\alpha:=\log_r s.
\]
Since \(r\neq s\) are primes, \(\alpha\notin \mathbb Z\). The divisors of \(n\) are exactly
\[
r^0,r^1,\dots,r^a,\qquad s r^0,s r^1,\dots,s r^a,
\]
that is,
\[
r^e \quad \text{with } e\in E:=\{0,1,\dots,a\}\cup \{\alpha,\alpha+1,\dots,\alpha+a\}.
\]
Thus the divisor set is the merge of \emph{two} geometric progressions with common ratio \(r\), or equivalently the exponent set is the union of \emph{two} arithmetic progressions with common difference \(1\).

For a fixed divisor \(d_i=r^{e_i}\), all later divisors \(d_j=r^{e_j}\) contribute
\[
\frac{1}{d_j-d_i}
=
\frac{1}{r^{e_i}}\cdot \frac{1}{r^{e_j-e_i}-1}.
\]
The key point is that for fixed \(i\), among the exponent differences \(e_j-e_i\) there is at most one element in \((0,1)\), and for each integer \(m\geq 1\) there are at most two elements in \([m,m+1)\). This makes the non-nearest-neighbour tail absolutely summable with a uniform constant.

This two-progressions phenomenon is the exact reason the argument works for \(r^a s\). It also isolates the first obstruction: once one has three or more progressions in exponent space (for instance \(r^a s_1 s_2\)), a fixed starting divisor can have \emph{more than one} subunit exponent difference, and the present charging argument to a single consecutive gap no longer closes.

\section{ATTACK PLAN (ROUND-2)}

The exact gap left after Round-1 was: produce any genuinely new unconditional infinite family for which the conjectured inequality can be proved with an absolute constant.

The claim to prove now is:

\medskip

\noindent\textbf{Claim.}
\emph{If \(n=r^a s\) with \(a\geq 1\) and distinct primes \(r,s\), then}
\[
S_{\mathrm{all}}(n)\leq 9(1+S_{\mathrm{gap}}(n)).
\]

\medskip

The plan is:

\begin{enumerate}
\item encode the divisor set as \(d_i=r^{e_i}\) with \(e_i\) in the two-progressions set \(E\);
\item for each fixed \(i\), bound the full tail \(\sum_{j>i}1/(d_j-d_i)\) by the nearest-neighbour term \(1/(d_{i+1}-d_i)\) plus a convergent error of size \(O(1/d_i)\);
\item sum over \(i\), and use the explicit bound on \(\sum_{d\mid n}1/d\).
\end{enumerate}

This overcomes the Round-1 obstacle because the tail is no longer treated by the generic harmonic-loss bound; instead it is controlled by the rigid two-progressions structure.

\section{WORK (ROUND-2)}

\subsection*{Step 1: exponent parametrisation}

Let
\[
E:=\{0,1,\dots,a\}\cup \{\alpha,\alpha+1,\dots,\alpha+a\},
\qquad \alpha:=\log_r s.
\]
Sort the elements of \(E\) in increasing order:
\[
e_1<e_2<\cdots<e_t,
\qquad t=2a+2.
\]
Then
\[
d_i=r^{e_i}\qquad (1\leq i\leq t)
\]
is exactly the increasing list of positive divisors of \(n=r^a s\).

For \(1\leq i<t\), define
\[
\Delta_i:=e_{i+1}-e_i>0.
\]
Then
\[
d_{i+1}-d_i
=
r^{e_i}(r^{\Delta_i}-1),
\]
so
\[
\frac{1}{d_{i+1}-d_i}
=
\frac{1}{r^{e_i}}\cdot \frac{1}{r^{\Delta_i}-1}.
\]

\subsection*{Step 2: a pointwise tail bound}

\paragraph{Lemma 1.}
For each fixed \(i\in\{1,\dots,t-1\}\),
\[
\sum_{j=i+1}^{t}\frac{1}{d_j-d_i}
\leq
\frac{1}{d_{i+1}-d_i}
+
\frac{2}{d_i}\sum_{m=1}^{\infty}\frac{1}{r^m-1}.
\]

\emph{Proof.}
Fix \(i<t\). Since \(d_j=r^{e_j}\), we have
\[
\sum_{j=i+1}^{t}\frac{1}{d_j-d_i}
=
\frac{1}{d_i}\sum_{j=i+1}^{t}\frac{1}{r^{e_j-e_i}-1}.
\tag{1}
\]

We now analyse the set of positive exponent differences
\[
D_i:=\{e_j-e_i:\ j=i+1,\dots,t\}.
\]

There are two cases for \(e_i\):

\begin{itemize}
\item If \(e_i\in\{0,1,\dots,a\}\), then the differences to exponents in the same progression are positive integers, while the differences to exponents in the other progression form a subset of an arithmetic progression
\[
\gamma+\mathbb Z
\]
for some non-integer real \(\gamma\) (namely \(\gamma\equiv \alpha \pmod{1}\)).

\item If \(e_i\in\{\alpha,\alpha+1,\dots,\alpha+a\}\), then the differences to exponents in the same progression are again positive integers, while the differences to exponents in the other progression form a subset of an arithmetic progression
\[
\gamma+\mathbb Z
\]
for some non-integer real \(\gamma\) (namely \(\gamma\equiv -\alpha \pmod{1}\)).
\end{itemize}

In either case:

\begin{enumerate}
\item there is at most one element of \(D_i\) in the interval \((0,1)\);
\item for each integer \(m\geq 1\), the interval \([m,m+1)\) contains at most one integer from \(D_i\) and at most one non-integer element from the translated progression, hence at most two elements of \(D_i\).
\end{enumerate}

Now \(\Delta_i=e_{i+1}-e_i\) is the smallest element of \(D_i\). Therefore
\[
\sum_{x\in D_i}\frac{1}{r^x-1}
\leq
\frac{1}{r^{\Delta_i}-1}
+
2\sum_{m=1}^{\infty}\frac{1}{r^m-1},
\]
because after removing the smallest term, every remaining \(x\in D_i\) lies in \([m,m+1)\) for some integer \(m\geq 1\), and there are at most two such \(x\) for each \(m\).

Substituting this into (1) gives
\[
\sum_{j=i+1}^{t}\frac{1}{d_j-d_i}
\leq
\frac{1}{d_i}\cdot \frac{1}{r^{\Delta_i}-1}
+
\frac{2}{d_i}\sum_{m=1}^{\infty}\frac{1}{r^m-1}
=
\frac{1}{d_{i+1}-d_i}
+
\frac{2}{d_i}\sum_{m=1}^{\infty}\frac{1}{r^m-1},
\]
as required.
\qquad \(\Box\)

\subsection*{Step 3: summing over all starting divisors}

Summing Lemma 1 over \(i=1,\dots,t-1\), we obtain
\[
S_{\mathrm{all}}(n)
=
\sum_{i=1}^{t-1}\sum_{j=i+1}^{t}\frac{1}{d_j-d_i}
\leq
\sum_{i=1}^{t-1}\frac{1}{d_{i+1}-d_i}
+
2\left(\sum_{m=1}^{\infty}\frac{1}{r^m-1}\right)\left(\sum_{i=1}^{t-1}\frac{1}{d_i}\right).
\]
Hence
\[
S_{\mathrm{all}}(n)
\leq
S_{\mathrm{gap}}(n)
+
2\left(\sum_{m=1}^{\infty}\frac{1}{r^m-1}\right)\left(\sum_{d\mid n}\frac{1}{d}\right).
\tag{2}
\]

We now estimate the two factors on the right-hand side.

\paragraph{Bound for the series.}
Since \(r\geq 2\),
\[
\sum_{m=1}^{\infty}\frac{1}{r^m-1}
\leq
\sum_{m=1}^{\infty}\frac{1}{2^m-1}.
\]
Also, for every \(m\geq 2\),
\[
2^m-1 \geq 3\cdot 2^{m-2},
\]
so
\[
\sum_{m=1}^{\infty}\frac{1}{2^m-1}
=
1+\sum_{m=2}^{\infty}\frac{1}{2^m-1}
\leq
1+\sum_{m=2}^{\infty}\frac{1}{3\cdot 2^{m-2}}
=
1+\frac{2}{3}
=
\frac{5}{3}.
\]
Therefore
\[
2\sum_{m=1}^{\infty}\frac{1}{r^m-1}
\leq
\frac{10}{3}.
\tag{3}
\]

\paragraph{Bound for the reciprocal divisor sum.}
Since the divisors of \(n=r^a s\) are \(r^u\) and \(sr^u\) for \(0\leq u\leq a\),
\[
\sum_{d\mid n}\frac{1}{d}
=
\left(1+\frac{1}{s}\right)\sum_{u=0}^{a}\frac{1}{r^u}.
\]
If \(r=2\), then \(s\geq 3\), so
\[
\left(1+\frac{1}{s}\right)\sum_{u=0}^{a}\frac{1}{r^u}
\leq
\frac{4}{3}\cdot 2
=
\frac{8}{3}.
\]
If \(r\geq 3\), then
\[
\sum_{u=0}^{a}\frac{1}{r^u}\leq \frac{1}{1-\frac{1}{3}}=\frac{3}{2},
\qquad
1+\frac{1}{s}\leq \frac{3}{2},
\]
hence
\[
\left(1+\frac{1}{s}\right)\sum_{u=0}^{a}\frac{1}{r^u}
\leq
\frac{9}{4}
<
\frac{8}{3}.
\]
Thus in all cases
\[
\sum_{d\mid n}\frac{1}{d}\leq \frac{8}{3}.
\tag{4}
\]

Substituting (3) and (4) into (2), we get
\[
S_{\mathrm{all}}(n)
\leq
S_{\mathrm{gap}}(n)
+
\frac{10}{3}\cdot \frac{8}{3}
=
S_{\mathrm{gap}}(n)+\frac{80}{9}.
\]
Since \(\frac{80}{9}<9\), it follows that
\[
S_{\mathrm{all}}(n)\leq S_{\mathrm{gap}}(n)+9 \leq 9(1+S_{\mathrm{gap}}(n)).
\]
This proves the theorem.
\qquad \(\Box\)

\subsection*{Consequences}

The Round-2 theorem gives an unconditional absolute-constant bound for every number of the form \(r^a s\) with \(r\neq s\) prime. Together with the Round-1 prime-power case, this verifies the conjectured type of estimate for the infinite families
\[
n=r^a
\qquad\text{and}\qquad
n=r^a s
\]
with distinct primes \(r,s\).

\section{ADVERSARIAL VERIFICATION}

I now check the new argument against failure modes.

\paragraph{1. Could \(\alpha=\log_r s\) be an integer?}
No. If \(\alpha\in\mathbb Z\), then \(s=r^\alpha\), impossible for distinct primes \(r\neq s\).

\paragraph{2. Are all divisors represented exactly once as \(r^{e_i}\)?}
Yes. The divisor set of \(r^a s\) is exactly
\[
\{r^u:0\leq u\leq a\}\cup \{sr^u:0\leq u\leq a\}
=
\{r^u:0\leq u\leq a\}\cup \{r^{\alpha+u}:0\leq u\leq a\}.
\]
These are distinct because \(\alpha\notin\mathbb Z\).

\paragraph{3. What if the two progressions do not interleave at all?}
That causes no problem. In that case there may be no exponent difference in \((0,1)\), but the proof of Lemma 1 still works: the smallest difference is then at least \(1\), and every interval \([m,m+1)\) still contains at most two relevant differences.

\paragraph{4. Did the proof accidentally use a hidden monotonicity or convexity statement?}
No. The only properties used are:
\begin{enumerate}
\item the exact representation of the exponents as a union of two arithmetic progressions with common step \(1\);
\item the monotonicity of \(x\mapsto 1/(r^x-1)\) for \(x>0\);
\item the explicit geometric-series bounds in (3) and (4).
\end{enumerate}

\paragraph{5. Does this conflict with the Round-1 conditional disproof?}
No. Tao's conditional counterexamples use products of \emph{many} nearly consecutive primes. The present theorem only covers the two-progressions case \(r^a s\), which is far outside that regime.

\paragraph{6. Where does this method first break?}
The first genuine barrier is the passage from two exponent progressions to three. If the divisor exponents are a union of three translates of \(\{0,1,\dots,a\}\), then from one starting exponent there can be more than one positive difference in \((0,1)\). The present argument charges only one exceptional near term to the consecutive-gap sum, so it no longer closes verbatim. This sharply isolates the first unresolved case for this method.

\section{FINAL}

\paragraph{\textbf{UNRESOLVED (BUT STRICTLY ADVANCED).}}

The original global Erd\H{o}s problem remains unresolved unconditionally. However, this round proves a new unconditional theorem strictly beyond Round-1:

\begin{quote}
For every \(n=r^a s\) with \(a\geq 1\) and \(r\neq s\) prime,
\[
S_{\mathrm{all}}(n)\leq 9\bigl(1+S_{\mathrm{gap}}(n)\bigr).
\]
\end{quote}

This is a substantive extension of the previously proved prime-power case and isolates a precise obstruction: the current method works exactly for divisor sets arising from a merge of two geometric progressions, and its first nontrivial failure mode is the three-progressions situation.

\section{COMPLETION ESTIMATE}

\noindent \textbf{COMPLETION: 82\%.}

\section{REFERENCES}

\begin{enumerate}
\item T. Tao, \emph{On the sum of reciprocals of gaps between divisors}, dated September 9, 2025.
\item T. F. Bloom, \emph{Erd\H{o}s Problem \#884}, \texttt{https://www.erdosproblems.com/884}, accessed 2026-03-07.
\end{enumerate}
